\section{Results}
\label{sec:results}

\subsection{What survived the withdrawal}
\label{sec:results:availability}

The audit reported in Table~\ref{tab:availability} establishes the scale of the loss. In
the goalkeeping detail table, \NKeeperAdvEmpty{} of \NKeeperAdvStatCols{} statistical
columns are entirely empty in both seasons, leaving three: nineties played, goals
conceded and penalties conceded. Neither of the two surviving substantive columns
supports any of the analyses that table exists to enable. The pattern is the same across
the outfield tables, where passing, possession, and goal and shot creation retain their
headers and lose their contents.

This is the first result of the paper, and it is a methodological one. A study that
selected its features from the published schema of this source would produce an empty
feature matrix while every column name resolved correctly, and the failure would surface
only as unexplained missingness far downstream. Feature selection against observed values
rather than published schemas is now a necessary step when working with public football
data.

\subsection{Correcting defensive counts for possession}
\label{sec:results:padj}

Before any structure is estimated, the defensive counts have to be separated from the
opportunity to record them. Across the \PadjTeamSeasons{} team-seasons available, team
possession ranges from \PadjPossessionMin{} to \PadjPossessionMax{} percent, and team
defensive volume does rise as opponent possession rises, so the confound is real. It is
much weaker than proportional, however. The fitted elasticities are
\ElasticityInterceptions{} for interceptions, \ElasticityTackles{} for tackles won and
\ElasticityFouls{} for fouls committed, against the value of one that the conventional
adjustment assumes.

That difference matters, and Table~\ref{tab:padj} shows why. The three raw correlations
with team possession are \PadjCorrRawInterceptions{}, \PadjCorrRawTackles{} and
\PadjCorrRawFouls{}, all negative, since a player whose team keeps the ball records fewer
defensive actions. Applying the conventional exponent overshoots every one of them: the
correlations become \PadjCorrUnitInterceptions{}, \PadjCorrUnitTackles{} and
\PadjCorrUnitFouls{}, reversing sign in all three cases and, for fouls committed, ending
further from zero than where it started. The correction passes through the confound and out
the other side. Adjusting with the measured elasticity instead leaves
\PadjCorrAdjInterceptions{}, \PadjCorrAdjTackles{} and \PadjCorrAdjFouls{}, which are
three to four times closer to zero than the conventional exponent achieves and are the
smallest residual dependence available from either choice.

Two qualifications belong with this. The correction is calibrated on a league whose
possession shares are relatively compressed, so the elasticities should not be transported
to a competition with a wider spread without refitting. More importantly, adjusting the
defensive counts does not make the feature set as a whole team-neutral, because the
attacking statistics still scale with a team's own share of the ball. What follows should
be read as having removed one identifiable confound, not all of them.

\begin{table}[t]
\centering
\small
\caption{Possession adjustment. The elasticity is fitted from team totals. The final three columns give the correlation between each per-90 rate and team possession: unadjusted, adjusted with the fitted elasticity, and adjusted with the conventional exponent of one. A value near zero is the goal.}
\label{tab:padj}
\begin{tabular}{lrrrr}
\toprule
 & & \multicolumn{3}{c}{Correlation with team possession} \\
\cmidrule(lr){3-5}
Statistic & Elasticity & Unadjusted & Adjusted & Exponent one \\
\midrule
Interceptions & 0.525 & -0.122 & -0.022 & 0.070 \\
Tackles won & 0.432 & -0.111 & -0.018 & 0.105 \\
Fouls committed & 0.283 & -0.102 & -0.037 & 0.129 \\
\bottomrule
\end{tabular}
\end{table}


\subsection{Why the correction is one sided}
\label{sec:results:symmetric}

Correcting only the defensive statistics invites an obvious objection. Attacking output
also depends on possession, so the feature set is still team dependent and the correction
looks like half a correction. The objection is worth testing rather than answering by
assertion, and the test produces the sharpest methodological result in the paper.

The two sides of the ball do not behave alike. Fitted the same way, against a team's own
possession rather than the opponent's, the attacking elasticities are \BetaShots{} for
shots, \BetaNPXG{} for non-penalty expected goals, \BetaGoals{} for non-penalty goals,
\BetaAssists{} for assists, \BetaXGChain{} for xGChain and \BetaXGBuildup{} for
xGBuildup. Every one of them except crossing, at \BetaCrosses{}, exceeds one, while every
defensive elasticity falls well below it. Figure~\ref{fig:elasticities} shows the two
families either side of the proportional line.

That asymmetry is not a curiosity, it is the whole argument. For defending, opponent
possession is an exposure denominator in the plain sense that a tackle cannot occur while
your own side holds the ball, and the sub-proportional response says exposure is exactly
what it measures. For attacking, an elasticity above one says the opposite: ten percent
more of the ball buys appreciably more than ten percent more output. Returns to possession
are increasing. That is a statement about football rather than about measurement, so
dividing it out would remove signal along with any bias. The two Understat involvement
measures make the point most starkly, since both exceed two: a side that keeps the ball
does not merely have more possessions, its players are involved in more of each.

The counterfactual was nonetheless computed. Applying the fitted exponents to the attacking
block as well moves the correlation between club centroids and final league position from
\SymCorrBaseline{} to \SymCorrAdjusted{}, so a symmetric correction does remove a real
part of the team relationship rather than nothing at all. It does not remove all of it, the
partition moves at an adjusted Rand index of \SymARI{} against the one reported here, and
the silhouette drifts from \SymSilhouetteBase{} to \SymSilhouetteAdj{}. The choice is
therefore a judgement rather than a fact, and it is recorded as one: we treat possession
driven productivity as something a role profile should reflect, not something to divide
away, and a reader who takes the opposite view has the numbers to see what that would cost.

\begin{figure}[t]
  \centering
  \includegraphics[width=\textwidth]{possession_elasticities.pdf}
  \caption{Possession elasticities for both families of statistic, fitted from team totals.
  Defensive volume responds to opponent possession less than proportionately, while
  attacking output responds to a team's own possession more than proportionately. Crossing
  is the one attacking statistic that behaves like a defensive one.}
  \label{fig:elasticities}
\end{figure}

\subsection{Descriptive structure}
\label{sec:results:descriptive}

After filtering, \NEligiblePrimary{} outfield players are eligible in the primary season,
comprising \NDFPrimary{} defenders, \NMFPrimary{} midfielders and \NFWPrimary{} forwards.
Summary statistics by position group appear in Table~\ref{tab:summary}, and the
distributions of four representative statistics are shown in
Figure~\ref{fig:distributions}.

The feature set is strongly non-normal. All \NFeaturesRejectNormal{} of the
\NFeatures{} features reject the Shapiro-Wilk null at the five percent level, mean
absolute skewness is \SkewMeanAbs{}, and the most skewed feature is goals per shot at
\SkewMax{}. Mardia's test rejects multivariate normality decisively on both the skewness
and kurtosis components. Principal component analysis requires no distributional
assumption to be computed, so this does not invalidate the decomposition, but it does
mean the components should be read as directions of variance rather than as latent
factors with any distributional interpretation, and it justifies the use of Spearman
rather than Pearson correlation throughout.

\begin{table}[t]
\centering
\small
\caption{Headline per-90 statistics by position group in the primary season, reported as mean with standard deviation in parentheses.}
\label{tab:summary}
\begin{tabular}{lrrr}
\toprule
Statistic & DF & MF & FW \\
\midrule
np\_xg\_p90 & 0.05 (0.04) & 0.16 (0.12) & 0.41 (0.17) \\
xa\_p90 & 0.06 (0.06) & 0.15 (0.10) & 0.16 (0.11) \\
key\_passes\_p90 & 0.51 (0.44) & 1.26 (0.63) & 1.05 (0.54) \\
xg\_chain\_p90 & 0.35 (0.16) & 0.50 (0.20) & 0.63 (0.20) \\
xg\_buildup\_p90 & 0.30 (0.14) & 0.28 (0.13) & 0.18 (0.09) \\
shots\_p90 & 0.54 (0.33) & 1.54 (0.85) & 2.53 (0.69) \\
shot\_accuracy\_pct & 29.22 (20.60) & 31.91 (13.77) & 40.30 (9.65) \\
crosses\_p90 & 1.22 (1.62) & 2.26 (1.75) & 1.40 (1.62) \\
interceptions\_p90 & 0.95 (0.37) & 0.74 (0.44) & 0.24 (0.16) \\
tackles\_won\_p90 & 1.06 (0.48) & 1.17 (0.55) & 0.55 (0.33) \\
fouls\_committed\_p90 & 0.87 (0.38) & 1.32 (0.55) & 1.27 (0.62) \\
offsides\_p90 & 0.06 (0.07) & 0.16 (0.17) & 0.53 (0.33) \\
\bottomrule
\end{tabular}
\end{table}


\begin{figure}[t]
  \centering
  \includegraphics[width=\textwidth]{distributions_by_position.pdf}
  \caption{Distributions of four representative per-90 statistics across position groups
  in the primary season. The separation between groups is clear for shot volume and for
  defensive actions, and much weaker for creation.}
  \label{fig:distributions}
\end{figure}

\subsection{Statistics that change meaning by position}
\label{sec:results:divergence}

The correlation analysis produced the clearest substantive finding available from the
reduced feature set. Correlations were computed within each position group and compared
across groups, and the ten pairs whose correlation varies most between groups are
reported in Table~\ref{tab:divergence} and shown in Figure~\ref{fig:divergence}. The
largest spread between position groups reaches \CorrSpreadMax{}, against a median spread
across all pairs of \CorrSpreadMedian{}.

Every one of the ten most divergent pairs changes sign between position groups. The
relationship between shots on target and crosses is positive for midfielders and negative
for forwards, which is to say that a midfielder who crosses often is generally a more
active attacker while a forward who crosses often is generally not the player getting on
the end of chances. The relationship between key passes and tackles won is positive for
defenders and negative for midfielders, distinguishing the full back whose defensive work
and creative output rise together from the midfield division of labour between a
destroyer and a creator.

The practical implication is that a single pooled correlation matrix, of the kind that
appears routinely in public football analysis, averages over these sign reversals and
reports a number that describes no actual population of players. It also supplies the
central justification for the position stratified analysis that follows, since a global
embedding fitted across all outfield players is fitted across relationships that point in
opposite directions.

\begin{table}[t]
\centering
\small
\caption{The ten feature pairs whose Spearman correlation differs most across position groups. Every pair changes sign between at least two groups.}
\label{tab:divergence}
\begin{tabular}{llrrrr}
\toprule
Feature A & Feature B & DF & MF & FW & Spread \\
\midrule
Shots on target & Crosses & -0.01 & 0.43 & -0.44 & 0.87 \\
Crosses & Shot accuracy \% & -0.22 & 0.26 & -0.57 & 0.82 \\
npxG & Crosses & -0.15 & 0.36 & -0.43 & 0.79 \\
Key passes & Tackles won & 0.45 & -0.34 & 0.23 & 0.79 \\
Crosses & Interceptions & 0.20 & -0.50 & 0.26 & 0.76 \\
Crosses & Tackles won & 0.50 & -0.25 & 0.47 & 0.75 \\
xA & Shots on target & 0.17 & 0.56 & -0.19 & 0.75 \\
Shots & Interceptions & 0.19 & -0.52 & -0.05 & 0.70 \\
xA & Tackles won & 0.39 & -0.32 & 0.11 & 0.70 \\
xA & Interceptions & 0.18 & -0.51 & 0.11 & 0.70 \\
\bottomrule
\end{tabular}
\end{table}


\begin{figure}[t]
  \centering
  \includegraphics[width=\textwidth]{corr_divergence_top10.pdf}
  \caption{The ten feature pairs whose within-group correlation differs most across
  position groups in the primary season. Every pair changes sign between at least two
  groups.}
  \label{fig:divergence}
\end{figure}

\subsection{Comparing embeddings}
\label{sec:results:embeddings}

Principal component analysis, UMAP and t-SNE were fitted to the same standardised
features and compared in Figure~\ref{fig:embeddings}. The selected UMAP configuration
used \UMAPBestNeighbors{} neighbours at a minimum distance of \UMAPBestMinDist{}, chosen
by trustworthiness over the twelve configuration grid reported in
Appendix~\ref{sec:appendix:umap}.

The three methods agree on the dominant structure and differ in how they present it.
Trustworthiness is \TrustPCA{} for the linear projection, \TrustUMAP{} for UMAP and
\TrustTSNE{} for t-SNE, while continuity is close to identical across the three. Nearest
neighbour position purity at ten neighbours is \PurityPCA{} for the linear projection and
\PurityUMAP{} for UMAP, so roughly two thirds to three quarters of a player's closest
statistical neighbours share his listed position group. That figure is informative in
both directions: it confirms that the surviving features carry real positional signal,
and it establishes that a substantial minority of every player's nearest neighbours are
listed at a different position, which is the space in which data-driven roles can differ
from listed ones.

All three embeddings display a continuous gradient from defenders through midfielders to
forwards rather than separated groups. That gradient, and not a set of islands, is the
honest description of the geometry, and it conditions how the clustering that follows
should be read.

\begin{figure}[t]
  \centering
  \includegraphics[width=\textwidth]{embedding_comparison.pdf}
  \caption{Two-dimensional embeddings of the same feature set, coloured by listed position
  group. The defender to forward gradient is continuous in all three, and t-SNE is shown
  for visual contrast only and is not used for clustering.}
  \label{fig:embeddings}
\end{figure}

\subsection{The third dimension}
\label{sec:results:threed}

Two components are a convenient way to draw the space and a poor way to describe it. The
first two carry shooting volume and wide creation and together account for
\ThreeDVarTwo{} of the variance. The third carries duels and discipline, loading on fouls
committed, cards and tackles won, and it is orthogonal to both of the others, so a planar
projection does not compress it, it deletes it. Retaining it raises captured variance to
\ThreeDVarThree{}. Figure~\ref{fig:threedpositions} shows the outfield space in all three,
drawn from two viewpoints so that structure can be told apart from the accident of a single
projection.

Adding the dimension does not change the headline geometry. The defender to forward gradient
remains continuous rather than resolving into islands, which is the same conclusion the
planar embeddings reached and is now established against a projection that had the
opportunity to falsify it. Figure~\ref{fig:threedclusters} places the selected partition in
the same space: the boundary cuts across the shooting axis and is close to indifferent to
duels and discipline, which is consistent with a split that is substantially about whether a
player shoots.

The third dimension does change the club picture, and this is the one place where it
produces a finding the planar analysis could not. Club centroids correlate with final
league position at \ThreeDTeamPCOne{} on the first component and \ThreeDTeamPCTwo{} on the
second, but only \ThreeDTeamPCThree{} on the third. Shooting volume and wide creation both
track league position; the physical dimension of duels, fouls and cards essentially does
not. A club does not finish higher by fouling more or less, which is unsurprising stated
plainly and is worth stating because the duelling axis is a full component of player
variation that carries almost no information about team success.
Figure~\ref{fig:threedteams} shows the clubs, labelled at the extremes of the table.

Finally, the three-dimensional view settles a structural question raised in
Section~\ref{sec:results:archetypes} by measurement rather than assertion. Drawing each
archetype centroid as a direction from the origin, the two centroids within a position
group have a cosine of \CosineDF{} for defenders, \CosineMF{} for midfielders and
\CosineFW{} for forwards. They are exactly anti-parallel, which is forced by standardising
within group before a two-way split. Figure~\ref{fig:threedaxes} shows the consequence: six
archetypes occupy three lines through the origin, not six independent directions.

\begin{figure}[t]
  \centering
  \includegraphics[width=\textwidth]{pca3d_positions.pdf}
  \caption{Outfield players in the first three principal components, from two viewpoints.
  Grey points on the floor are the vertical projection of each player, supplying a depth
  cue that a static rendering otherwise lacks.}
  \label{fig:threedpositions}
\end{figure}

\begin{figure}[t]
  \centering
  \includegraphics[width=\textwidth]{pca3d_clusters.pdf}
  \caption{The selected global partition in three components, with cluster centroids marked
  by crosses. The split runs along the shooting axis and is close to indifferent to the
  third component.}
  \label{fig:threedclusters}
\end{figure}

\begin{figure}[t]
  \centering
  \includegraphics[width=\textwidth]{pca3d_teams.pdf}
  \caption{Clubs as minutes-weighted centroids in three components, shaded by final league
  position with darker shading for a higher finish. Vertical lines to the floor give the
  depth cue, and only the top and bottom of the table are labelled.}
  \label{fig:threedteams}
\end{figure}

\begin{figure}[t]
  \centering
  \includegraphics[width=\textwidth]{pca3d_archetype_axes.pdf}
  \caption{Archetype centroid directions, drawn as unit vectors so the figure shows
  orientation rather than magnitude. Each position group forms a single line through the
  origin, so the six archetypes describe three axes.}
  \label{fig:threedaxes}
\end{figure}

\subsection{How many roles the data supports}
\label{sec:results:clustering}

Applying the consensus rule declared in Section~\ref{sec:methods} to the four selection
criteria gives $k=\KGlobal{}$ for the global clustering and the same value within each of
the three position groups. The full criterion curves appear in
Figure~\ref{fig:kselection} and Appendix~\ref{sec:appendix:kselection}. This is the second
principal result of the paper. The surviving public feature set supports a binary
division within each position group and nothing finer. Requesting more clusters produces
partitions that the criteria do not support and that resampling does not reproduce.

Stability separates the position groups sharply. Across \BootstrapN{} bootstrap resamples
the mean adjusted Rand index against the reference partition is \BootARIGlobal{} globally
and \BootARIMF{} for midfielders, both of which describe a partition that reappears almost
unchanged under resampling. For defenders the mean falls to \BootARIDF{} with a standard
deviation of \BootARISDDF{}, and for forwards to \BootARIFW{} with a standard deviation of
\BootARISDFW{}. The minimum adjusted Rand index observed across forward resamples is
\BootARIMinFW{}, which is to say that in the least favourable resample the partition
recovered is no better than chance. Forwards number only \NFWPrimary{} players above the
minutes threshold, so this is partly a sample size problem, but it is a real limit on what
can be claimed. The consensus co-assignment structure is shown in
Figure~\ref{fig:consensus}.

The relationship between discovered clusters and listed positions is weak. The global
partition scores an adjusted Rand index of \ARIvsPosition{} against the four-way position
group with a cluster purity of \PurityVsPosition{}, so the clusters are not a rediscovery
of the positions that the source already records. Figure~\ref{fig:clustervsposition} shows
where the two disagree. Read together with the embedding geometry, the interpretation is
that the dominant axis in the surviving data is not defender against forward but
attacking involvement against defensive involvement, which cuts across the listed groups
and places a ball-playing full back nearer an ordinary midfielder than a defensive
counterpart at his own listed position.

\begin{figure}[t]
  \centering
  \includegraphics[width=\textwidth]{k_selection_curves.pdf}
  \caption{Cluster selection criteria across candidate cluster counts. The consensus rule
  was fixed before fitting and selects the same value in every scope.}
  \label{fig:kselection}
\end{figure}

\begin{figure}[t]
  \centering
  \begin{subfigure}[t]{0.48\textwidth}
    \includegraphics[width=\textwidth]{consensus_matrix.pdf}
    \caption{Bootstrap co-assignment.}
    \label{fig:consensus}
  \end{subfigure}
  \hfill
  \begin{subfigure}[t]{0.48\textwidth}
    \includegraphics[width=\textwidth]{cluster_vs_position.pdf}
    \caption{Clusters against listed positions.}
    \label{fig:clustervsposition}
  \end{subfigure}
  \caption{Cluster stability and the correspondence with listed positions.}
  \label{fig:clustervalidation}
\end{figure}

\begin{figure}[t]
  \centering
  \includegraphics[width=0.75\textwidth]{cluster_scatter_pca.pdf}
  \caption{The selected partition shown in the principal component space.}
  \label{fig:clusterscatter}
\end{figure}

\subsection{Goalkeepers}
\label{sec:results:keepers}

The goalkeeper analysis is the clearest illustration of what the withdrawal cost.
\NKeeperAdvEmpty{} of the \NKeeperAdvStatCols{} statistical columns in the advanced
goalkeeping table are empty, and the \NKeeperAdvUsable{} that survive are duplicated in
the basic table, so the advanced table contributes nothing whatsoever. Of the five
composite axes the study set out to build, \KeeperAxesBuilt{} can be constructed and
\KeeperAxesLost{} cannot. No proxy is substituted for the missing four, because no
surviving column measures sweeping, distribution or cross handling even approximately.

With post-shot expected goals unavailable, shot stopping is estimated by regressing saves
on shots on target faced and treating the residual as saves above volume. The regression
returns a coefficient of determination of \KeeperSavesRSq{}, which is high almost by
construction, since every shot on target is either saved or conceded. The informative
quantity is the residual spread, a standard error of \KeeperSavesResidSE{} saves across
the season. This index adjusts for how much work a goalkeeper receives but not for how
hard that work is, which is precisely what post-shot expected goals provided, so it should
be read as a weak substitute rather than an equivalent.

Save percentage shows no relationship with workload. The correlation is
\KeeperWorkloadR{} in the primary season at a p-value of \KeeperWorkloadP{}, and it
changes sign to \KeeperWorkloadRRepl{} in the replication season. A relationship whose
sign is not stable across two seasons is the strongest available form of a null result.

Goalkeeper clustering does not produce meaningful structure, and we report that rather
than presenting a partition. The same consensus rule selects \KeeperKPrimary{} clusters in
the primary season at a silhouette of \KeeperSilhouettePrimary{} but \KeeperKRepl{}
clusters in the replication season at \KeeperSilhouetteRepl{}, so the cluster count does
not replicate. The smallest replication cluster contains \KeeperMinClusterRepl{} members,
below the floor fixed before fitting. Most decisively, cluster membership explains more
variance in situational quantities than in goalkeeping ones: the proportion of variance
explained is \KeeperEtaTeamPoints{} for team points per match against \KeeperEtaSavePct{}
for save percentage. The partition recovers the strength of the team in front of the
goalkeeper, not the goalkeeper. No archetype names are assigned.

\begin{figure}[t]
  \centering
  \begin{subfigure}[t]{0.48\textwidth}
    \includegraphics[width=\textwidth]{keeper_saves_vs_sota.pdf}
    \caption{Saves against shots on target faced.}
    \label{fig:keepersaves}
  \end{subfigure}
  \hfill
  \begin{subfigure}[t]{0.48\textwidth}
    \includegraphics[width=\textwidth]{keeper_savepct_vs_workload.pdf}
    \caption{Save percentage against workload.}
    \label{fig:keeperworkload}
  \end{subfigure}
  \caption{Goalkeeper shot stopping and workload, the only two dimensions the surviving
  data supports.}
  \label{fig:keepers}
\end{figure}

Figure~\ref{fig:keeperranking} ranks eligible goalkeepers on the saves above volume
index, and Figure~\ref{fig:keepercomposites} places them on the two axes that survive.
The distribution of both composites is given in
Figure~\ref{fig:keepercompositedistributions}, and the correlation structure of the
surviving goalkeeper features in Figure~\ref{fig:keepercorr}. The partition itself is
shown in Figure~\ref{fig:keeperclusters}, presented for completeness rather than as a
finding, for the reasons given above.

\begin{figure}[t]
  \centering
  \includegraphics[width=\textwidth]{keeper_shot_stopping_ranking.pdf}
  \caption{Goalkeepers ranked by saves above the volume of shots on target faced. The
  index adjusts for how much work a goalkeeper receives but not for its difficulty.}
  \label{fig:keeperranking}
\end{figure}

\begin{figure}[t]
  \centering
  \begin{subfigure}[t]{0.48\textwidth}
    \includegraphics[width=\textwidth]{keeper_composites.pdf}
    \caption{The two surviving composite axes.}
    \label{fig:keepercomposites}
  \end{subfigure}
  \hfill
  \begin{subfigure}[t]{0.48\textwidth}
    \includegraphics[width=\textwidth]{keeper_clusters_pca.pdf}
    \caption{The goalkeeper partition, shown for completeness.}
    \label{fig:keeperclusters}
  \end{subfigure}
  \caption{Goalkeeper composites and the partition that the validation checks reject.}
  \label{fig:keeperpanel}
\end{figure}

\subsection{Archetypes}
\label{sec:results:archetypes}

Because the selection rule returned two clusters in each position group, the study
describes six outfield archetypes rather than the richer taxonomy the design anticipated.
Their centroid profiles are shown in Figure~\ref{fig:archetypeprofiles} and their radar
comparisons in Figure~\ref{fig:radarpanels}. Names were assigned only after inspecting the
distinguishing features and the member players, and the justification for each is
reproduced from the generated record rather than composed for the paper.

Among defenders the split separates \NDFZero{} players with no positive attacking
signature at all from \NDFOne{} whose key passes, assists, expected assists and crosses
all sit around one standard deviation above the defensive mean. The first group is named
for what it does not do, because every one of its eight distinguishing features is
negative and its largest positive deviation is negligible. The second corresponds closely
to full backs and wing backs. Among midfielders the division is the cleanest in the study,
separating \NMFZero{} players whose only positive deviations are defensive or disciplinary
from \NMFOne{} whose shooting and creation measures are all elevated. Among forwards the
split distinguishes \NFWZero{} central players who shoot from \NFWOne{} wide players who
cross and create.

One structural property must be stated because it limits how much the six profiles can be
read as independent. Features are standardised within position group before clustering, so
with two clusters the two centroids in each group are necessarily anti-parallel: their
size-weighted sum is zero. Each position group therefore yields two ends of a single axis
rather than two independently shaped profiles, which is why the two members of each pair
mirror one another so exactly. The axis is real and the ordering along it is informative,
but the six archetypes represent three axes, not six independent types.

Outlier analysis is reported in Figure~\ref{fig:outliers}. The isolation forest and the
Mahalanobis distance agree closely, which is reassuring given that one is a partition
method and the other a distance method.

\begin{figure}[t]
  \centering
  \includegraphics[width=\textwidth]{archetype_profiles.pdf}
  \caption{Centroid deviation of each archetype from its position group mean, across the
  full feature set. Within each position group the two profiles are necessarily mirror
  images, since the features are standardised within group before clustering.}
  \label{fig:archetypeprofiles}
\end{figure}

\begin{figure}[t]
  \centering
  \begin{subfigure}[t]{0.32\textwidth}
    \includegraphics[width=\textwidth]{radar_panel_DF.pdf}
    \caption{Defenders.}
    \label{fig:radardf}
  \end{subfigure}
  \hfill
  \begin{subfigure}[t]{0.32\textwidth}
    \includegraphics[width=\textwidth]{radar_panel_MF.pdf}
    \caption{Midfielders.}
    \label{fig:radarmf}
  \end{subfigure}
  \hfill
  \begin{subfigure}[t]{0.32\textwidth}
    \includegraphics[width=\textwidth]{radar_panel_FW.pdf}
    \caption{Forwards.}
    \label{fig:radarfw}
  \end{subfigure}
  \caption{Archetype radars on fixed axes within each position group.}
  \label{fig:radarpanels}
\end{figure}

\subsection{What a supervised model sees}
\label{sec:results:supervised}

The clustering results leave one question open: are the surviving features too weak to
carry positional information at all, or do they carry it while the clustering chooses to
partition along a different axis? A supervised model settles it. Predicting the listed
position group from the same \NFeatures{} features, gradient boosting reaches an accuracy
of \LGBMAccuracy{} and a macro F1 of \LGBMMacroF{}, and logistic regression reaches
\LogitAccuracy{} and \LogitMacroF{}, against a majority class baseline of
\BaselineAccuracy{} and \BaselineMacroF{}. The features clearly do carry position, so the
weak agreement between clusters and positions reflects the clustering finding a different
and equally real axis, not an empty feature set. The confusion is almost entirely between
adjacent groups, as Figure~\ref{fig:confusion} shows, and forwards are the weakest class.

The SHAP decomposition in Figures~\ref{fig:shapbeeswarm} and \ref{fig:shapbyclass} carries
a finding that follows directly from the withdrawal. With possession and passing volume
gone, positions are identified largely by absence. Defenders are recognised by low
shooting and low key passes, forwards by low interceptions, and midfielders by an
intermediate profile on both. A classifier of this kind, trained before January 2026,
would have identified a defender by what he did with the ball. It now identifies him by
what he does not do.

Players the classifier places in the wrong group overlap significantly with the players
the clustering flagged as hybrids. The overlap rate is \OverlapLogit{} for logistic
regression at an enrichment of \EnrichmentLogit{} over independence, and
\OverlapLGBM{} at an enrichment of \EnrichmentLGBM{} for gradient boosting. The
enrichment is modest because the hybrid flag is itself broad, covering \NHybridPlayers{}
of \NEligiblePrimary{} players, so this is corroboration rather than a strong result.

The finishing above role regression illustrates a hazard worth recording. Predicting
non-penalty expected goals per ninety from non-shooting behaviour returns a cross
validated coefficient of determination of \FinishingRSq{}. An earlier specification that
also included xGChain returned \FinishingRSqLeaky{}, which looks far better and is in fact
circular: Understat credits xGChain for possessions the player finished himself, so the
difference between xGChain and xGBuildup contains the target and correlates with it at
\FinishingLeakCorr{}. The apparently strong model was reconstructing the outcome rather
than predicting it. Even at the honest value, roughly half the variance is unexplained, so
the residual leaderboard in Figure~\ref{fig:finishing} is presented as an illustration and
not as a measurement of finishing ability.

\begin{figure}[t]
  \centering
  \begin{subfigure}[t]{0.48\textwidth}
    \includegraphics[width=\textwidth]{confusion_matrix.pdf}
    \caption{Confusion between listed positions.}
    \label{fig:confusion}
  \end{subfigure}
  \hfill
  \begin{subfigure}[t]{0.48\textwidth}
    \includegraphics[width=\textwidth]{shap_by_class.pdf}
    \caption{Feature importance per position.}
    \label{fig:shapbyclass}
  \end{subfigure}
  \caption{Supervised validation of the surviving feature set.}
  \label{fig:supervised}
\end{figure}

\subsection{Team signatures}
\label{sec:results:teams}

Aggregating players into clubs produces the one result in this study that is both strong
and stable. The minutes-weighted position of a club's outfield players on the first
principal component correlates with final league position at a Spearman coefficient of
\TeamPCOneRho{} in the primary season and \TeamPCOneRhoRepl{} in the replication season,
which is close to an exact reproduction. Since the first component is essentially
attacking involvement, this says that clubs whose players collectively shoot and create
more finish higher, which is unsurprising in itself, but the near identical coefficient
across two independent seasons shows the team level summary is measuring something stable
even where the player level partitions are not.

Hierarchical clustering on archetype composition groups the \NTeams{} clubs into
\NStyleGroups{} style families, shown in Figure~\ref{fig:dendrogram}, with composition by
club in Figure~\ref{fig:composition}. Figure~\ref{fig:teampanels} places each club's
players in the shared embedding against the rest of the league, which makes squad shape
legible in a way that a composition bar cannot: clubs differ visibly in how far their
defenders sit toward the midfield region.

\begin{figure}[t]
  \centering
  \begin{subfigure}[t]{0.48\textwidth}
    \includegraphics[width=\textwidth]{team_dendrogram.pdf}
    \caption{Style families.}
    \label{fig:dendrogram}
  \end{subfigure}
  \hfill
  \begin{subfigure}[t]{0.48\textwidth}
    \includegraphics[width=\textwidth]{team_composition.pdf}
    \caption{Archetype composition by club.}
    \label{fig:composition}
  \end{subfigure}
  \caption{Club stylistic signatures, ordered by final league position.}
  \label{fig:teamsignatures}
\end{figure}

\begin{figure}[p]
  \centering
  \includegraphics[width=\textwidth]{team_panels.pdf}
  \caption{Each club's outfield players highlighted in the shared embedding, with the rest
  of the league in grey. Axes are shared across all twenty panels.}
  \label{fig:teampanels}
\end{figure}

\subsection{Similarity, replication and noise}
\label{sec:results:novel}

\subsubsection{Player similarity}
Nearest neighbour search within position group recovers football sense for players whose
roles are defined by the surviving statistics, and fails informatively for players whose
roles are not. Full backs return full backs, penalty box forwards return penalty box
forwards, and centre backs return centre backs. The instructive failures are the
midfielders. A holding midfielder's nearest neighbours are creative attacking midfielders,
because with no tackles by pitch third, no blocks, no clearances, no pressures and no
passing volume, the only things left that distinguish him are his set piece deliveries and
key passes. The metric finds the players who share that signature and misses the entire
defensive half of his job. An inverted full back returns three centre backs for the same
reason, since the role is defined by possession behaviour that no longer exists in the
data. Similarity scores among centre backs are systematically the lowest in the table,
which is the same feature poverty appearing as an inability to tell defenders apart at
all. Figure~\ref{fig:similarity} shows the structure.

\subsubsection{Cross-season replication}
\NCommonPlayers{} players are eligible in both seasons. The declared rule independently
selects the same number of clusters in every scope in both seasons, so the structure
replicates. Membership replicates only partly. The adjusted Rand index between seasons is
\CrossSeasonARIGlobal{} globally, with \CrossSeasonChangedGlobal{} players changing
cluster, and within position group it falls to \CrossSeasonARIDF{} for defenders,
\CrossSeasonARIMF{} for midfielders and \CrossSeasonARIFW{} for forwards. The within group
figures are worse than the bootstrap stability of Section~\ref{sec:results:clustering},
which is what would be expected if the within group splits are finer than the data can
reliably support. Figure~\ref{fig:crossseason} shows the transitions.

\subsubsection{Sensitivity to sampling noise}
\label{sec:results:shrinkage}
Empirical Bayes shrinkage toward the position group mean, with the degree of shrinkage set
by minutes played, provides the sharpest test of whether the partitions describe playing
roles or sampling noise. Shrinkage is substantial rather than cosmetic, and the results
divide cleanly.

The global partition survives. Only \ShrinkChangedGlobal{} of \NEligiblePrimary{} players
change cluster, the adjusted Rand index between the original and shrunken partitions is
\ShrinkARIGlobal{}, the selected number of clusters is unchanged, and the rate at which
players move falls monotonically with minutes played, exactly as the noise account
predicts.

The within group partitions fare worse, and they do not fare equally. For defenders the
adjusted Rand index falls to \ShrinkARIDF{}, with \ShrinkChangedDF{} of \NDFPrimary{}
players moving, which is a partition largely dissolved by removing sampling noise.
Midfielders retain \ShrinkARIMF{} and forwards \ShrinkARIFW{}, both of which survive in
part rather than collapsing. The forward figure is notably higher than it was before the
possession correction of Section~\ref{sec:results:padj} was applied, which is itself
informative: some of what previously looked like noise in the forward partition was
exposure to opponent possession, and removing it left a more reproducible split. The mechanism is
visible in the estimated variance components: several features are assigned an estimated
between-player variance of zero within a position group, meaning the model attributes
their entire observed spread to sampling noise and collapses them to the group mean,
removing them from the clustering altogether. Shot accuracy retains very little
reliability, because its true spread between players is small relative to the noise in
estimating it from a season of shots.

Read alongside the bootstrap and replication results, the three tests do not deliver a
single uniform verdict, and reporting them as though they did would overstate the case.
Defenders fail all three: bootstrap stability of \BootARIDF{}, cross-season agreement of
\CrossSeasonARIDF{}, and near dissolution under shrinkage. Forwards fail two of the three,
with a bootstrap minimum of \BootARIMinFW{} across resamples and cross-season agreement of
\CrossSeasonARIFW{}, while partly surviving shrinkage. Midfielders pass all three. The
conclusion that survives every test is therefore narrower than a single summary sentence
would suggest: the global division between attacking and defensive involvement is a real
property of players, the midfield split is trustworthy, and the defender split should not
be used to characterise an individual at all.
Figures~\ref{fig:shrinkage} and \ref{fig:shrinkageminutes} show the effect.

\begin{figure}[t]
  \centering
  \begin{subfigure}[t]{0.48\textwidth}
    \includegraphics[width=\textwidth]{similarity_heatmap.pdf}
    \caption{Nearest neighbour similarity.}
    \label{fig:similarity}
  \end{subfigure}
  \hfill
  \begin{subfigure}[t]{0.48\textwidth}
    \includegraphics[width=\textwidth]{cross_season_transitions.pdf}
    \caption{Between-season transitions.}
    \label{fig:crossseason}
  \end{subfigure}
  \caption{Similarity structure and cross-season replication.}
  \label{fig:novelpanel}
\end{figure}

\subsubsection{Role drift within a season}
\label{sec:results:drift}
Every result so far treats a season as one observation per player, which assumes a
player's profile is homogeneous across it. Match reports make that assumption testable,
and because the withdrawal removed the affected columns from the match schema rather than
emptying them, what remains there is fully populated. A full season is 380 reports at seven
seconds each, so the scrape is restricted to one club, \DriftTeam{}, chosen because its
eligible players span all six archetypes and therefore give drift the widest range to move
across. \DriftMatches{} matches were retrieved and \DriftPlayers{} players accumulated at
least one complete rolling window of \DriftWindow{} matches.

The match feature set has \DriftFeatures{} of the eighteen features used elsewhere, since
the Understat expected goals block has no match equivalent. Rather than assume that the
coarser space is a fair stand-in, we refit the season clustering using only those
\DriftFeatures{} features and compare it to the partition reported throughout. The two
agree at an adjusted Rand index of \DriftReducedARI{}, so the reduced space reproduces
almost all of the main division and drift measured within it is interpretable.

Most players do not drift. \DriftSwitchers{} of \DriftPlayers{}, or \DriftSwitchShare{}
percent, change cluster at least once during the season, and the remainder occupy a single
cluster throughout. The players who stay put are the ones whose roles are least ambiguous,
holding midfielders on one side and the recognised attackers on the other, while the movers
are predominantly full backs and wide players whose attacking involvement varies with the
fixture. Figure~\ref{fig:drift} shows the trajectories.

The finding cuts both ways and should be read that way. For four fifths of players a season
aggregate describes a single role and the homogeneity assumption holds. For the remaining
fifth it averages over two, which means their season-level archetype is a summary of
something that changed rather than a description of something stable. That is a limitation
of season aggregates in general rather than of this feature set in particular, and it is one
that only match level data can expose.

\begin{figure}[t]
  \centering
  \includegraphics[width=\textwidth]{role_drift.pdf}
  \caption{Rolling cluster assignment across the season for the \DriftTeam{} players who
  change cluster at least once, on a \DriftWindow{} match window. Players not shown occupied
  one cluster for the whole season.}
  \label{fig:drift}
\end{figure}

\subsubsection{Age and archetype}
\label{sec:results:age}
The two archetypes within a position group differ mainly in attacking involvement, and
involvement is plausibly something a career passes through rather than a fixed trait, so
the association with age is worth testing. It is tested within position group, because
archetypes are nested inside positions and a pooled table would largely recover the fact
that forwards and defenders have different age profiles.

Only defenders show a relationship. The crossing and chance creating defenders average
\AgeGapDF{} years younger than the low involvement group, a difference that the rank test
places at \AgeKruskalPDF{}, and the share of the crossing archetype falls monotonically
from the youngest band to the oldest. Midfielders show nothing at all, with a gap of
\AgeGapMF{} years at \AgeKruskalPMF{}, and the global partition is unrelated to age at
\AgeGlobalP{}. Forwards point the same way as defenders but cannot support the claim,
since six of the ten cells in their table have expected counts below five.

The defender result also illustrates a point about test choice that is easy to get wrong.
A chi-square on the banded table returns \AgeChiPDF{} and would ordinarily be read as no
association. That test spends its degrees of freedom on any departure from independence,
while the pattern here is monotonic, and a rank correlation against age recovers it at
\AgeTrendRhoDF{} with a p-value of \AgeTrendPDF{}. The banded test is not wrong, it is
underpowered against the alternative that actually holds, and reporting only the chi-square
would have converted a real finding into a null one. Figure~\ref{fig:age} shows the shares.

\begin{figure}[t]
  \centering
  \includegraphics[width=\textwidth]{age_archetype.pdf}
  \caption{Share of each age band taken by each archetype, within position group. The
  defender panel falls monotonically with age, which a chi-square on bands does not detect
  and a rank correlation does.}
  \label{fig:age}
\end{figure}

\begin{figure}[t]
  \centering
  \begin{subfigure}[t]{0.48\textwidth}
    \includegraphics[width=\textwidth]{shrinkage_effect.pdf}
    \caption{Effect of shrinkage on assignments.}
    \label{fig:shrinkage}
  \end{subfigure}
  \hfill
  \begin{subfigure}[t]{0.48\textwidth}
    \includegraphics[width=\textwidth]{shrinkage_by_minutes.pdf}
    \caption{Assignment changes by minutes played.}
    \label{fig:shrinkageminutes}
  \end{subfigure}
  \caption{Empirical Bayes shrinkage as a test of whether the partitions describe roles or
  noise.}
  \label{fig:shrinkagepanel}
\end{figure}

\begin{figure}[t]
  \centering
  \begin{subfigure}[t]{0.48\textwidth}
    \includegraphics[width=\textwidth]{outliers_scatter.pdf}
    \caption{Outliers within position group.}
    \label{fig:outliers}
  \end{subfigure}
  \hfill
  \begin{subfigure}[t]{0.48\textwidth}
    \includegraphics[width=\textwidth]{finishing_above_role.pdf}
    \caption{Residuals from the role model.}
    \label{fig:finishing}
  \end{subfigure}
  \caption{Unusual players, and the residual leaderboard that should not be read as a
  measurement of finishing ability.}
  \label{fig:outlierspanel}
\end{figure}

\begin{figure}[t]
  \centering
  \includegraphics[width=0.85\textwidth]{shap_beeswarm.pdf}
  \caption{Global SHAP values for the gradient boosted position classifier. Shooting
  volume and interceptions carry most of the signal, so a player's position is now
  identified largely by the actions he does not perform.}
  \label{fig:shapbeeswarm}
\end{figure}
